\chapter{Teorema Thales}\label{ch:g9:thales}

Bagaimana kamu mengukur tinggi sebuah piramida tanpa memanjatnya?
Thales dari Miletus membandingkan bayangannya dengan bayangan sebuah
tongkat. Teoremanya --- \omterm{def:g6:lines:perp}{garis sejajar} memotong \omterm{def:g6:lines:objects}{ruas garis} menjadi
bagian yang \omterm{def:g9:linfunc:linear}{sebanding} --- adalah jantung matematika model berskala,
peta, dan perbesaran, sekaligus salah satu teorema bernama yang
tertua.

\section{Teoremanya}

\begin{theorem}[Thales]\label{thm:g9:thales:direct}
Misalkan dua garis bertemu di sebuah titik $A$. Ambillah $M$ pada
garis pertama dan $N$ pada garis kedua, dengan $B$ lebih jauh sepanjang
garis pertama dan $C$ lebih jauh sepanjang garis kedua. Kalau garis
$(MN)$ dan $(BC)$ \emph{\omterm{def:g6:lines:perp}{sejajar}}, maka segitiga $AMN$ dan $ABC$ punya
sisi yang \omterm{def:g9:linfunc:linear}{sebanding}:
\[
\frac{AM}{AB} = \frac{AN}{AC} = \frac{MN}{BC}.
\]
\end{theorem}
\admitted

\begin{omfigure}
\begin{tikzpicture}[scale=1.05]
  \coordinate (A) at (0,0);
  \coordinate (B) at (5.0,0.6);
  \coordinate (C) at (3.2,2.8);
  \coordinate (M) at ($(A)!0.55!(B)$);
  \coordinate (N) at ($(A)!0.55!(C)$);
  \draw[thick] (A) -- (B) -- (C) -- cycle;
  \draw[omThm, very thick] (M) -- (N);
  \draw[omThm, very thick] (B) -- (C);
  \fill (A) circle (1.5pt) node[left, font=\small] {$A$};
  \fill (B) circle (1.5pt) node[right, font=\small] {$B$};
  \fill (C) circle (1.5pt) node[above, font=\small] {$C$};
  \fill (M) circle (1.5pt) node[below, font=\small] {$M$};
  \fill (N) circle (1.5pt) node[above left, font=\small] {$N$};
\end{tikzpicture}\hspace{1.1cm}%
\begin{tikzpicture}[scale=1.05]
  \coordinate (A) at (0,0);
  \coordinate (B) at (2.6,-1.3);
  \coordinate (C) at (1.4,-2.2);
  \coordinate (M) at ($(A)!-0.62!(B)$);
  \coordinate (N) at ($(A)!-0.62!(C)$);
  \draw[thick] (M) -- (B);
  \draw[thick] (N) -- (C);
  \draw[omThm, very thick] (M) -- (N);
  \draw[omThm, very thick] (B) -- (C);
  \fill (A) circle (1.5pt) node[right=2pt, font=\small] {$A$};
  \fill (B) circle (1.5pt) node[right, font=\small] {$B$};
  \fill (C) circle (1.5pt) node[below, font=\small] {$C$};
  \fill (M) circle (1.5pt) node[left, font=\small] {$M$};
  \fill (N) circle (1.5pt) node[above, font=\small] {$N$};
\end{tikzpicture}

{\small Kedua susunan Thales: segitiga bersarang (kiri) dan susunan
``kupu-kupu'', tempat $M$ dan $N$ duduk di sisi lain $A$ (kanan). Pada
keduanya, $(MN) \parallel (BC)$ dan ketiga perbandingannya sama.}
\end{omfigure}

\begin{method}[Menghitung sebuah panjang dengan Thales]\label{met:g9:thales:compute}
\begin{enumerate}
  \item Kenalilah titik $A$ tempat kedua garisnya berpotongan dan
        kedua \omterm{def:g6:lines:perp}{garis sejajarnya}; namailah kedua segitiganya;
  \item tulislah ketiga perbandingan yang sama, selalu ``segitiga
        kecil dibagi segitiga besar'';
  \item tahanlah kesamaan kedua perbandingan yang memuat ketiga
        panjang yang diketahui dan panjang yang belum diketahui;
  \item selesaikan perbandingan yang dihasilkannya dengan perkalian
        silang.
\end{enumerate}
\end{method}

\begin{example}\label{ex:g9:thales:compute}
Pada susunan kiri di atas, andaikan $AM = 3$, $AB = 5$, $AN = 2.4$
dan $MN = 3.3$, lalu marilah kita hitung $AC$ dan $BC$. Thales memberi
\[
\frac{AM}{AB} = \frac{AN}{AC} = \frac{MN}{BC},
\qquad\text{yaitu}\qquad
\frac35 = \frac{2.4}{AC} = \frac{3.3}{BC}.
\]
Dari $\frac35 = \frac{2.4}{AC}$, dengan mengalikan silang:
$3 \times AC = 5 \times 2.4 = 12$, jadi $AC = 4$. Dari
$\frac35 = \frac{3.3}{BC}$: $3 \times BC = 16.5$, jadi $BC = 5.5$.
\end{example}

\begin{example}[Thales dalam kehidupan nyata]\label{ex:g9:thales:stick}
Sebuah tongkat tegak setinggi $1$ m melemparkan bayangan $1.5$ m,
sedangkan sebuah pohon melemparkan bayangan $12$ m pada saat yang
sama. Sinar matahari \omterm{def:g6:lines:perp}{sejajar}, jadi segitiga tongkat-bayangan dan
segitiga pohon-bayangan berada pada susunan Thales:
\[
\frac{\text{tinggi pohonnya}}{1} = \frac{12}{1.5},
\qquad\text{jadi pohonnya setinggi } 8 \text{ m.}
\]
\end{example}

\section{Konversnya}

\begin{theorem}[Konvers Thales]\label{thm:g9:thales:converse}
Dengan titiknya diletakkan seperti pada \cref{thm:g9:thales:direct}
($M$ pada $(AB)$, $N$ pada $(AC)$, dengan urutan yang sama pada kedua
garisnya): kalau
\[
\frac{AM}{AB} = \frac{AN}{AC},
\]
maka garis $(MN)$ dan $(BC)$ \omterm{def:g6:lines:perp}{sejajar}.
\end{theorem}
\admitted

\begin{method}[Membuktikan atau menyangkal kesejajaran]\label{met:g9:thales:converse}
\begin{enumerate}
  \item Hitunglah secara terpisah kedua perbandingannya
        $\dfrac{AM}{AB}$ dan $\dfrac{AN}{AC}$ (sebagai \omterm{def:g9:fractions:fraction}{pecahan}, bukan
        sebagai \omterm{def:g6:decimals:rounding}{pembulatan});
  \item kalau keduanya sama \emph{dan} titiknya berurutan sama pada
        kedua garisnya, maka garisnya \omterm{def:g6:lines:perp}{sejajar} (menurut konvers Thales);
  \item kalau keduanya berbeda, garisnya \emph{tidak} \omterm{def:g6:lines:perp}{sejajar} ---
        karena kalau \omterm{def:g6:lines:perp}{sejajar}, teorema Thales akan memaksa
        perbandingannya sama.
\end{enumerate}
\end{method}

\begin{example}\label{ex:g9:thales:converse}
Pada garis 1, $AM = 4$ dan $AB = 10$; pada garis 2, $AN = 6$ dan
$AC = 15$. Maka
\[
\frac{AM}{AB} = \frac{4}{10} = \frac25,
\qquad
\frac{AN}{AC} = \frac{6}{15} = \frac25 :
\]
perbandingannya sama, urutannya sama: jadi $(MN) \parallel (BC)$.

Dengan $AC = 14$ sebagai gantinya: $\frac{6}{14} = \frac37 \neq
\frac25$, jadi garisnya tidak \omterm{def:g6:lines:perp}{sejajar}.
\end{example}

\section{Perbesaran dan pengecilan}

\begin{definition}[Menskalakan sebuah bangun]\label{def:g9:thales:scaling}
Memperbesar atau mengecilkan sebuah bangun dengan \emph{faktor
skala}\index{faktor skala} $k > 0$ berarti mengalikan \emph{semua}
panjangnya dengan $k$: disebut \emph{perbesaran} ketika $k > 1$, dan
\emph{pengecilan} ketika $k < 1$. Pada susunan Thales, segitiga $AMN$
adalah pengecilan $ABC$ dengan faktor $k = \frac{AM}{AB}$.
\end{definition}

\begin{proposition}[Pengaruhnya pada sudut dan luas]\label{prop:g9:thales:scaling}
Menskalakan dengan faktor $k$ mempertahankan sudut dan bentuk
bangunnya, mengalikan setiap panjang dengan $k$, dan mengalikan setiap
\emph{luas} dengan $k^2$.
\end{proposition}
\admitted

\begin{omfigure}
\begin{tikzpicture}[scale=0.72]
  \draw[thick, omDef, fill=omDef!8] (0,0) rectangle (2,1.2);
  \draw[thick, omThm, fill=omThm!8] (3.4,0) rectangle (7.4,2.4);
  \node[font=\small, omDef] at (1,-0.45) {$2 \times 1.2$};
  \node[font=\small, omThm] at (5.4,-0.45) {$4 \times 2.4$};
  \node[font=\small] at (1,0.6) {luas $2.4$};
  \node[font=\small] at (5.4,1.2) {luas $9.6$};
\end{tikzpicture}

{\small Melipatduakan panjangnya ($k = 2$) mengalikan luasnya dengan
$k^2 = 4$: empat salinan persegi panjang kecilnya mengubin yang besar.}
\end{omfigure}

\begin{example}
Sebuah foto $10 \times 15$ cm diperbesar dengan \omterm{def:g9:thales:scaling}{faktor skala} $k = 3$:
cetakannya berukuran $30 \times 45$ cm. Luasnya berubah dari $150$
cm$^2$ menjadi $150 \times 9 = 1350$ cm$^2$ --- yaitu $9$ kali lebih
besar, bukan $3$.
\end{example}

\section{Latihan}

\begin{exercise}[$\star$]\label{exo:g9:thales:1}
Garis $(MN)$ dan $(BC)$ \omterm{def:g6:lines:perp}{sejajar}, dengan $M$ pada $[AB]$ dan $N$ pada
$[AC]$; $AM = 2$, $AB = 6$, $AN = 3$, $BC = 9$. Hitunglah $AC$ dan $MN$.
\end{exercise}

\begin{exercise}[$\star$]\label{exo:g9:thales:2}
Susunan yang sama: $AM = 5$, $MB = 3$ (hati-hati: $MB$, bukan $AB$!),
$AN = 4$. Hitunglah $AC$.
\end{exercise}

\begin{exercise}[$\star$]\label{exo:g9:thales:3}
Pada susunan kupu-kupu, $A$ terletak antara $M$ dan $B$ dan antara $N$
dan $C$, dengan $(MN) \parallel (BC)$; $AM = 3$, $AB = 7.5$, $AN = 2$,
$MN = 2.6$. Hitunglah $AC$ dan $BC$.
\end{exercise}

\begin{exercise}[$\star$]\label{exo:g9:thales:4}
$M$ berada pada $[AB]$ dengan $AM = 6$, $AB = 8$; $N$ berada pada
$[AC]$ dengan $AN = 9$, $AC = 12$. Apakah garis $(MN)$ dan $(BC)$ \omterm{def:g6:lines:perp}{sejajar}?
\end{exercise}

\begin{exercise}[$\star\star$]\label{exo:g9:thales:5}
Sama seperti di atas dengan $AM = 4$, $AB = 6$, $AN = 5$, $AC = 8$.
Apakah $(MN)$ dan $(BC)$ \omterm{def:g6:lines:perp}{sejajar}? Berilah alasan dengan hati-hati.
\end{exercise}

\begin{exercise}[$\star\star$]\label{exo:g9:thales:6}
Seseorang setinggi $1.8$ m berdiri $2$ m dari kaki sebuah lampu jalan;
bayangannya berukuran $3$ m. Gambarlah susunan Thales yang dibentuk
lampunya, orangnya, dan ujung bayangannya, lalu hitunglah tinggi
lampunya.
\end{exercise}

\begin{exercise}[$\star\star$]\label{exo:g9:thales:7}
Sebuah peta punya skala $1 : 25\,000$ (panjang pada petanya adalah
panjang yang sebenarnya dikalikan $k = \frac{1}{25000}$).
\begin{enumerate}
  \item Dua desa berjarak $6.8$ cm pada petanya. Berapa jaraknya yang
        sebenarnya, dalam km?
  \item Sebuah hutan berluas $8$ cm$^2$ pada petanya. Berapa luasnya
        yang sebenarnya, dalam km$^2$?
\end{enumerate}
\end{exercise}

\begin{exercise}[$\star\star$]\label{exo:g9:thales:8}
Segitiga $ABC$ punya $AB = 12$, $AC = 15$, $BC = 18$. Titik $M$ pada
$[AB]$ memenuhi $AM = 8$, dan garis melalui $M$ yang \omterm{def:g6:lines:perp}{sejajar} dengan
$(BC)$ memotong $[AC]$ di $N$.
\begin{enumerate}
  \item Hitunglah $AN$ dan $MN$.
  \item Berapa \omterm{def:g9:thales:scaling}{faktor skala} dari $ABC$ ke $AMN$? Bandingkan
        \omterm{def:g6:measure:perimeter}{keliling} kedua segitiganya.
\end{enumerate}
\end{exercise}

\begin{exercise}[$\star\star\star$]\label{exo:g9:thales:9}
Misalkan $ABCD$ sebuah trapesium dengan $(AB) \parallel (CD)$,
$AB = 4$ dan $CD = 6$, yang diagonalnya $[AC]$ dan $[BD]$ bertemu di
$O$. Dengan memakai Thales pada susunan kupu-kupu di sekitar $O$,
hitunglah perbandingan $\dfrac{OA}{OC}$, lalu tunjukkan bahwa $O$
memotong kedua diagonalnya dengan perbandingan yang sama.
\end{exercise}

\section{Soal: Penggaris tanpa ukuran, kamera lubang jarum, dan
ketiga rata-rata sebuah trapesium}

\begin{problem}[{Soal akhir pekan --- Thales yang bekerja: membagi
ruas garis mana pun menjadi bagian yang sama, mengukur Matahari
dengan kotak sepatu, dan trapesium tempat ketiga rata-rata terkenalnya
bertemu}]
\label{pb:g9:thales:1}
Teorema Thales adalah matematika \emph{sinar}: sinar matahari,
sinar cahaya lewat lubang jarum, sinar pensil dari sebuah \omterm{def:g5:solids:def}{titik sudut}.
Soal ini memakainya dengan tiga cara --- sebagai alat pelukisan
(membagi \omterm{def:g6:lines:objects}{ruas garis} \emph{sepanjang apa pun}, bahkan $\sqrt2$,
menjadi bagian yang benar-benar sama), sebagai alat ukur (yaitu \omterm{def:g4:geometry:circle}{diameter}
Matahari, dengan kotak sepatu dan sekeping uang logam), dan sebagai
kaca pembesar pada trapesium
\cref{exo:g9:thales:9}, tempat rata-rata aritmetika, geometri dan
harmonik pada seri ini semuanya muncul dalam satu gambar.

\medskip
\noindent\textbf{Bagian I --- Membagi dengan penggaris tanpa ukuran.}
\begin{enumerate}
  \item Gambarlah \omterm{def:g6:lines:objects}{ruas garis} $[AB]$ sepanjang $7$ cm. Untuk
        memotongnya menjadi tiga bagian yang sama dengan jangka dan
        penggaris tanpa ukuran: gambarlah sinar apa saja dari $A$
        (yang tidak melalui $B$); langkahkan tiga panjang jangka yang
        sama padanya, yang memberi titik $P_1$, $P_2$, $P_3$;
        hubungkan $P_3$ ke $B$; lalu gambarlah \omterm{def:g6:lines:perp}{garis sejajar}
        $(P_3 B)$ yang melalui $P_1$ dan $P_2$. Lakukan pelukisannya
        lalu tandai tempat \omterm{def:g6:lines:perp}{garis sejajarnya} memotong $[AB]$.
  \item Berilah alasan dengan \cref{thm:g9:thales:direct} bahwa kedua
        titik yang ditandai itu memotong $[AB]$ tepat di titik
        sepertiganya.
  \item Sesuaikan caranya untuk membagi sebuah \omterm{def:g6:lines:objects}{ruas garis} menjadi $5$
        bagian yang sama, lalu untuk melukis $\frac35$ \omterm{def:g6:lines:objects}{ruas garis} yang
        diberikan.
  \item Lukislah titik $M$ pada $[AB]$ dengan
        $\frac{AM}{MB} = \frac23$ (yaitu $M$ pada $\frac25$ jalan
        dari $A$). Berapa langkah yang sama pada sinarnya, dan
        titik yang mana yang dihubungkan ke $B$?
  \item Mengapa pelukisan ini lebih baik daripada mengukur dengan
        penggaris berskala? Tinjaulah \omterm{def:g6:lines:objects}{ruas garis} sepanjang $\sqrt2$
        (yaitu diagonal persegi satuan, yang irasional menurut
        \cref{pb:g9:sqrt:1}): apa yang akan diberikan pengukuran, dan
        apa yang diberikan Thales?
\end{enumerate}

\medskip
\noindent\textbf{Bagian II --- Kamera lubang jarum.}
Tusuklah sebuah jarum menembus satu sisi kotak yang tertutup: pada
sisi yang di seberangnya, muncul bayangan dunia yang terbalik. Sebuah
titik benda, lubang jarumnya, dan titik bayangannya segaris --- karena
cahaya merambat lurus --- jadi bendanya (tinggi $H$, pada jarak $D$
di depan lubangnya) dan bayangannya (tinggi $h$, pada dinding
belakang sedalam $d$ di belakang lubangnya) duduk pada susunan
kupu-kupu milik Thales.
\begin{enumerate}[resume]
  \item Turunkan rumus lubang jarumnya
        $h = H \times \frac dD$ dari Thales pada susunan kupu-kupu
        di sekitar lubangnya.
  \item Sebuah pohon $12$ m berdiri $30$ m dari kotak sepatu sedalam
        $20$ cm. Berapa tinggi bayangannya, dan menghadap ke mana?
  \item Arahkan kotaknya ke Matahari: pada $d = 1$ m di belakang
        lubang jarumnya, bayangan Mataharinya adalah cakram
        berdiameter kira-kira $9$ mm. Diberikan jarak Mataharinya
        $D = 1.5 \times 10^{11}$ m, hitunglah \omterm{def:g4:geometry:circle}{diameternya} dalam
        \omterm{def:g9:fractions:scientific}{notasi ilmiah} (yaitu notasi pada \cref{pb:g9:fractions:1}
        yang sedang bekerja). Bandingkan dengan nilai sebenarnya,
        $1.39 \times 10^9$ m.
  \item Bulan: \omterm{def:g4:geometry:circle}{diameternya} $3.5 \times 10^6$ m, jaraknya
        $3.8 \times 10^8$ m. Hitunglah perbandingan
        $\frac{\text{diameter}}{\text{jarak}}$ untuk Bulan
        dan untuk Matahari. Kebetulan beruntung apa yang disingkap
        kedua bilangannya --- dan peristiwa menakjubkan apa yang
        dimungkinkannya?
  \item Dalam satu atau dua kalimat: mengapa bayangan lubang
        jarumnya terbalik, dan mengapa memperbesar lubang jarumnya
        membuat bayangannya lebih terang tetapi lebih kabur?
\end{enumerate}

\medskip
\noindent\textbf{Bagian III --- Trapesium ketiga rata-ratanya.}
$ABCD$ adalah trapesium pada \cref{exo:g9:thales:9}:
$(AB) \parallel (CD)$, $AB = 4$, $CD = 6$, dengan diagonal yang
berpotongan di $O$ dan $\frac{OA}{OC} = \frac{OB}{OD} = \frac{4}{6} =
\frac23$.
\begin{enumerate}[resume]
  \item Latihan kupu-kupu: dua garis berpotongan di $O$ dengan
        $(MN) \parallel (PQ)$ pada kedudukan kupu-kupu; $OM = 3$,
        $OP = 7.5$, $MN = 4$ dan $OQ = 6$. Hitunglah $PQ$ dan
        $ON$.
  \item Jelaskan bagaimana Thales, yang dipakai dengan perbandingan
        $\frac12$, memuat teorema titik tengah pada
        \cref{thm:g8:midpoints:direct} sebagai keadaan khususnya.
  \item Gambarlah \omterm{def:g6:lines:perp}{garis sejajar} alasnya yang melalui $O$; garis itu
        memotong $[AD]$ di $E$. Dengan bekerja pada segitiga $ACD$
        (perhatikan $\frac{AO}{AC} = \frac25$), hitunglah $EO$.
  \item Menurut perhitungan cerminnya pada segitiga $BDC$, potongan
        $OF$ (dengan $F$ pada $[BC]$) punya panjang yang sama.
        Simpulkan bahwa
        \[
        EF = 2 \times \frac{4 \times 6}{4 + 6} = 4.8 :
        \]
        \omterm{def:g6:lines:perp}{garis sejajar} yang melalui potongan diagonalnya berukuran
        \emph{\omterm{pb:g8:speed:1}{rata-rata harmonik}} kedua alasnya --- yaitu rata-rata
        perjalanan pulang pergi pada \cref{pb:g8:speed:1}, yang
        muncul kembali dalam geometri murni.
  \item Penutup akbarnya: hitunglah ketiga rata-rata klasik
        kedua alasnya $4$ dan $6$ --- aritmetika
        $\frac{4+6}{2}$, geometri $\sqrt{4 \times 6}$, harmonik
        $\frac{2 \times 4 \times 6}{4 + 6}$ --- lalu urutkan.
        Masing-masingnya hidup dalam trapesiumnya: rata-rata
        aritmetikanya adalah garis tengah yang menghubungkan titik
        tengah kedua kakinya (\cref{pb:g7:areas:1}), \omterm{pb:g8:speed:1}{rata-rata
        harmoniknya} adalah \omterm{def:g6:lines:objects}{ruas garis} $EF$ milikmu, dan \omterm{pb:g8:cosine:1}{rata-rata
        geometrinya} adalah \omterm{def:g6:lines:perp}{garis sejajar}
        yang memotong trapesiumnya menjadi dua trapesium yang
        \emph{sebangun} (periksalah nilainya dengan kalkulator;
        buktinya tantangan tambahan yang manis). Di mana
        ketaksamaan antara ketiganya sudah dibuktikan dalam buku ini
        (\cref{pb:g8:cosine:1}, \cref{pb:g8:speed:1})?

\end{enumerate}
\end{problem}
